% !TEX root = ../../ctfp-print.tex

\lettrine[lhang=0.17]{到}{目前为止，我们主要}处理单个范畴或一对范畴的情况。在某些情况下这种限制显得过于严格。

例如，在定义范畴$\cat{C}$中的极限时，我们引入了索引范畴$\cat{I}$作为形成锥体基础的模板。此时本应再引入另一个范畴——一个平凡范畴——作为锥顶点的模板。但我们却使用了常值函子$\Delta_c$从$\cat{I}$到$\cat{C}$。

现在是时候修正这个不协调之处了。让我们通过三个范畴来定义极限。首先考虑从索引范畴$\cat{I}$到$\cat{C}$的函子$D$。这个函子选择了锥的基础——即图示函子。

\begin{figure}[H]
  \centering
  \includegraphics[width=0.4\textwidth]{images/kan2.jpg}
\end{figure}

新增的是包含单个对象（和单个恒等态射）的范畴$\cat{1}$。存在唯一可能的函子$K$从$\cat{I}$到该范畴。它将所有对象映射到$\cat{1}$中的唯一对象，所有态射映射到恒等态射。任何从$\cat{1}$到$\cat{C}$的函子$F$都为我们选择一个潜在的锥顶点。

\begin{figure}[H]
  \centering
  \includegraphics[width=0.4\textwidth]{images/kan15.jpg}
\end{figure}

锥是一个自然变换$\varepsilon$从$F \circ K$到$D$。注意$F \circ K$的作用与原来的$\Delta_c$完全相同。下图展示了这个变换：

\begin{figure}[H]
  \centering
  \includegraphics[width=0.4\textwidth]{images/kan3-e1492120491591.jpg}
\end{figure}

现在我们可以定义一个泛性质来选取"最佳"的这种函子$F$。这个$F$会将$\cat{1}$映射到$\cat{C}$中$D$的极限对象，而自然变换$\varepsilon$从$F \circ K$到$D$将提供相应的投影。这个泛函子称为$D$沿$K$的右Kan扩展（right Kan extension），记作$\Ran_{K}D$。

让我们表述这个泛性质。假设存在另一个锥——即另一个函子$F'$和自然变换$\varepsilon'$从$F' \circ K$到$D$：

\begin{figure}[H]
  \centering
  \includegraphics[width=0.4\textwidth]{images/kan31-e1492120512209.jpg}
\end{figure}

如果Kan扩展$F = \Ran_{K}D$存在，则必然存在唯一的自然变换$\sigma$从$F'$到它，使得$\varepsilon'$通过$\varepsilon$分解：
$$\varepsilon' = \varepsilon\ .\ (\sigma \circ K)$$
这里$\sigma \circ K$是两个自然变换的水平合成（其中一个是$K$上的恒等自然变换）。然后这个变换再与$\varepsilon$进行垂直合成。

\begin{figure}[H]
  \centering
  \includegraphics[width=0.4\textwidth]{images/kan5.jpg}
\end{figure}

在分量形式下，当作用于$\cat{I}$中的对象$i$时，我们得到：
$$\varepsilon'_i = \varepsilon_i \circ \sigma_{K i}$$
在我们的案例中，$\sigma$只有一个分量对应$\cat{1}$的唯一对象。因此，这确实是从由$F'$定义的锥顶点到由$\Ran_{K}D$定义的泛锥顶点的唯一态射。交换条件正好满足极限定义的要求。

但重要的是，我们现在可以自由地将平凡范畴$\cat{1}$替换为任意范畴$\cat{A}$，右Kan扩展的定义依然成立。

\section{右Kan扩展}

函子$D \Colon \cat{I} \to \cat{C}$沿函子$K \Colon \cat{I} \to \cat{A}$的右Kan扩展是一个函子$F \Colon \cat{A} \to \cat{C}$（记作$\Ran_{K}D$）连同自然变换
$$\varepsilon \Colon F \circ K \to D$$
使得对于任意其他函子$F' \Colon \cat{A} \to \cat{C}$和自然变换
$$\varepsilon' \Colon F' \circ K \to D$$
存在唯一的自然变换
$$\sigma \Colon F' \to F$$
使得$\varepsilon'$分解：
$$\varepsilon' = \varepsilon\ .\ (\sigma \circ K)$$
这看起来相当复杂，但可以通过下面这个优美的图示直观理解：

\begin{figure}[H]
  \centering
  \includegraphics[width=0.4\textwidth]{images/kan7.jpg}
\end{figure}

有趣的是，我们可以注意到某种意义上Kan扩展就像"函子乘法"的逆运算。有些作者甚至使用记号$D/K$表示$\Ran_{K}D$。在这种记号下，$\varepsilon$的定义（也称为右Kan扩展的余单位）看起来就像简单的约简：
$$\varepsilon \Colon D/K \circ K \to D$$
Kan扩展还有另一种解释。考虑函子$K$将范畴$\cat{I}$嵌入到$\cat{A}$内部。在最简单的情形下$\cat{I}$可以只是$\cat{A}$的一个子范畴。我们有一个函子$D$将$\cat{I}$映射到$\cat{C}$。能否将$D$扩展为在整个$\cat{A}$上定义的函子$F$？理想情况下，这种扩展会使复合函子$F \circ K$与$D$同构。换句话说，$F$将把$D$的定义域扩展到$\cat{A}$。但通常很难达到这么强的要求，我们只需满足一半条件即可——即存在从$F \circ K$到$D$的单向自然变换$\varepsilon$。（左Kan扩展则选择相反方向）

\begin{figure}[H]
  \centering
  \includegraphics[width=0.4\textwidth]{images/kan6.jpg}
\end{figure}

当然，当函子$K$在对象上不是单射或在hom集上不是忠实映射时（如极限的例子），这种嵌入图景就会失效。在这种情况下，Kan扩展会尽力推断丢失的信息。

\section{作为伴随的Kan扩展}

现在假设对于任意$D$（以及固定的$K$）右Kan扩展都存在。在这种情况下，$\Ran_{K}-$（短横线代替$D$）是从函子范畴${[}\cat{I}, \cat{C}{]}$到函子范畴${[}\cat{A}, \cat{C}{]}$的函子。事实证明，该函子正是预合成函子$- \circ K$的右伴随。后者将${[}\cat{A}, \cat{C}{]}$中的函子映射到${[}\cat{I}, \cat{C}{]}$中的函子。这个伴随关系表示为：
$$[\cat{I}, \cat{C}](F' \circ K, D) \cong [\cat{A}, \cat{C}](F', \Ran_{K}D)$$
这本质上是对我们称为$\varepsilon'$的每个自然变换对应唯一一个我们称为$\sigma$的自然变换这一事实的重新表述。

\begin{figure}[H]
  \centering
  \includegraphics[width=0.4\textwidth]{images/kan92.jpg}
\end{figure}

\noindent
进一步地，若我们选择范畴$\cat{I}$与$\cat{C}$相同，则可用恒等函子$I_{\cat{C}}$替代$D$。此时得到如下恒等式：
$$[\cat{C}, \cat{C}](F' \circ K, I_{\cat{C}}) \cong [\cat{A}, \cat{C}](F', \Ran_{K}I_{\cat{C}})$$
我们可以令$F'$等于$\Ran_{K}I_{\cat{C}}$。在这种情况下，右侧包含恒等自然变换，对应地左侧给出如下自然变换：
$$\varepsilon \Colon \Ran_{K}I_{\cat{C}} \circ K \to I_{\cat{C}}$$
这非常类似于伴随的上单位：
$$\Ran_{K}I_{\cat{C}} \dashv K$$
事实上，恒等函子沿函子$K$的右Kan扩展可用于计算$K$的左伴随。为此需要满足一个额外条件：右Kan扩展必须被函子$K$保持。所谓保持扩展是指，当我们计算与$K$预合成的函子的Kan扩展时，结果应与对原始Kan扩展与$K$预合成的结果一致。在本例中，该条件简化为：
$$K \circ \Ran_{K}I_{\cat{C}} \cong \Ran_{K}K$$
注意到使用除法记号$/$时，伴随可写作：
$$I/K \dashv K$$
这验证了我们的直觉——伴随描述某种逆元性质。此时保持条件变为：
$$K \circ I/K \cong K/K$$
函子沿自身的右Kan扩展$K/K$被称为余密度单子（codensity monad）。

伴随公式是一个重要结论，因为正如我们即将看到的，我们可以通过端（共端）来计算Kan扩展，从而获得寻找右（和左）伴随的实际方法。

\section{左Kan扩张}

有一种对偶构造给出了我们的左Kan扩张。为了建立一些直观理解，我们可以从余极限(colimit)的定义开始，通过重构将其与单元素范畴$\cat{1}$结合使用。我们利用函子$D \Colon \cat{I} \to \cat{C}$构建余锥(base)，并通过函子$F \Colon \cat{1} \to \cat{C}$选取其顶点(apex)。

\begin{figure}[H]
  \centering
  \includegraphics[width=0.4\textwidth]{images/kan81.jpg}
\end{figure}

\noindent
余锥的侧面即注入映射(injections)，它们构成了从$D$到$F \circ K$的自然变换$\eta$的分量。

\begin{figure}[H]
  \centering
  \includegraphics[width=0.4\textwidth]{images/kan10a.jpg}
\end{figure}

\noindent
余极限是泛余锥(universal cocone)。因此对于任意其他函子$F'$及自然变换
$$\eta' \Colon D \to F' \circ K$$

\begin{figure}[H]
  \centering
  \includegraphics[width=0.4\textwidth]{images/kan10b.jpg}
\end{figure}

\noindent
存在唯一的自然变换$\sigma$从$F$到$F'$

\begin{figure}[H]
  \centering
  \includegraphics[width=0.4\textwidth]{images/kan14.jpg}
\end{figure}

\noindent
使得：
$$\eta' = (\sigma \circ K)\ .\ \eta$$
这可通过下图示说明：

\begin{figure}[H]
  \centering
  \includegraphics[width=0.4\textwidth]{images/kan112.jpg}
\end{figure}

\noindent
将单元素范畴$\cat{1}$替换为$\cat{A}$后，该定义自然推广到左Kan扩张(left Kan extension)的概念，记作$\Lan_{K}D$。

\begin{figure}[H]
  \centering
  \includegraphics[width=0.4\textwidth]{images/kan12.jpg}
\end{figure}

\noindent
自然变换：
$$\eta \Colon D \to \Lan_{K}D \circ K$$
称为左Kan扩张的单位元(unit)。

与之前类似，我们可以将自然变换之间的一一对应：
$$\eta' = (\sigma \circ K)\ .\ \eta$$
用伴随关系(adjunction)重新表述：
$$[\cat{A}, \cat{C}](\Lan_{K}D, F') \cong [\cat{I}, \cat{C}](D, F' \circ K)$$
换句话说，左Kan扩张是预复合函子(precomposition functor)的左伴随(left adjoint)，而右Kan扩张则是其右伴随(right adjoint)。

正如恒等函子的右Kan扩张可用于计算$K$的左伴随，恒等函子的左Kan扩张实际上成为$K$的右伴随（其中$\eta$为伴随关系的单位元）：
$$K \dashv \Lan_{K}I_{\cat{C}}$$
结合这两个结果可得：
$$\Ran_{K}I_{\cat{C}} \dashv K \dashv \Lan_{K}I_{\cat{C}}$$

\section{作为端的Kan扩张}

Kan扩张的真正威力在于它们可以通过端（end）和余端（coend）进行计算。为了简化问题，我们将注意力限制在目标范畴$\cat{C}$为$\Set$的情形，但相关公式可推广至任意范畴。

让我们重新审视Kan扩张能够将函子作用扩展到原始定义域之外的思想。假设$K$将$\cat{I}$嵌入到$\cat{A}$内部。函子$D$将$\cat{I}$映射到$\Set$。我们可以说对于$K$像中的任意对象$a=Ki$，扩展的函子将$a$映射到$Di$。问题是该如何处理那些不在$K$像中的$\cat{A}$对象？关键思想是：每个这样的对象通过大量态射与$K$像中的每个对象潜在连接。函子必须保持这些态射关系。从对象$a$到$K$像的全体态射由Hom函子表征：
$$\cat{A}(a, K -)$$

\begin{figure}[H]
  \centering
  \includegraphics[width=0.4\textwidth]{images/kan13.jpg}
\end{figure}

\noindent
注意这个Hom函子是两个函子的复合：
$$\cat{A}(a, K -) = \cat{A}(a, -) \circ K$$
右Kan扩张是函子复合的右伴随：
$$[\cat{I}, \Set](F' \circ K, D) \cong [\cat{A}, \Set](F', \Ran_{K}D)$$
当我们将$F'$替换为Hom函子时：
$$[\cat{I}, \Set](\cat{A}(a, -) \circ K, D) \cong [\cat{A}, \Set](\cat{A}(a, -), \Ran_{K}D)$$
代入复合结构得到：
$$[\cat{I}, \Set](\cat{A}(a, K -), D) \cong [\cat{A}, \Set](\cat{A}(a, -), \Ran_{K}D)$$
右侧可通过Yoneda引理化简：
$$[\cat{I}, \Set](\cat{A}(a, K -), D) \cong \Ran_{K}D a$$
现在我们可以将自然变换集合写成端的形式，得到右Kan扩张的便捷公式：
$$\Ran_{K}D a \cong \int_i \Set(\cat{A}(a, K i), D i)$$
左Kan扩张则有对应的余端公式：
$$\Lan_{K}D a = \int^i \cat{A}(K i, a)\times{}D i$$
为验证其确实构成函子复合的左伴随：
$$[\cat{A}, \Set](\Lan_{K}D, F') \cong [\cat{I}, \Set](D, F' \circ K)$$
将公式代入左侧：
$$[\cat{A}, \Set](\int^i \cat{A}(K i, -)\times{}D i, F')$$
这是自然变换集合，可改写为端：
$$\int_a \Set(\int^i \cat{A}(K i, a)\times{}D i, F' a)$$
利用Hom函子的连续性，可用端代替余端：
$$\int_a \int_i \Set(\cat{A}(K i, a)\times{}D i, F' a)$$
应用乘积-指数伴随：
$$\int_a \int_i \Set(\cat{A}(K i, a),\ (F' a)^{D i})$$
指数对象同构于相应Hom集合：
$$\int_a \int_i \Set(\cat{A}(K i, a),\ \Set(D i, F' a))$$
根据Fubini定理允许交换两端顺序：
$$\int_i \int_a \Set(\cat{A}(K i, a),\ \Set(D i, F' a))$$
内层端代表两个函子间的自然变换集合，应用Yoneda引理得：
$$\int_i \Set(D i, F' (K i))$$
这正是我们要证明的伴随右侧自然变换集合：
$$[\cat{I}, \Set](D, F' \circ K)$$
这种使用端、余端和Yoneda引理的计算方式，正是端的"微积分"典型特征。

\section{Haskell中的Kan扩张}

Kan扩张的端/余端公式可以很容易地转换到Haskell中。我们从右扩张开始：
$$\Ran_{K}D a \cong \int_i \Set(\cat{A}(a, K i), D i)$$
我们将端替换为全称量词，将hom-集替换为函数类型：

\src{snippet01}
观察这个定义，显然\Ran必须包含一个类型为\a的值用于应用函数，以及两个函子\code{k}和\code{d}之间的自然变换。例如，假设\code{k}是树函子，\code{d}是列表函子，并且你获得了一个\code{Ran Tree {[}{]} String}。如果你传入以下函数：

\src{snippet02}
你会得到一个\code{Int}列表，依此类推。右Kan扩张会使用你的函数生成一棵树，然后将其重新打包成列表。例如，你可以传入一个从字符串生成解析树的解析器，得到的结果将是该树的深度优先遍历对应的列表。

通过将函子\d替换为恒等函子，右Kan扩张可以用来计算给定函子的左伴随。这导致函子\k的左伴随表示为以下多态函数集合：

\src{snippet03}
假设\code{k}是从幺半群范畴导出的遗忘函子。此时全称量词遍历所有幺半群。当然，在Haskell中我们无法表达幺半群定律，但以下是对结果自由函子的合理近似（遗忘函子\code{k}在对象上是恒等的）：

\src{snippet04}
正如预期，它生成自由幺半群，即Haskell列表：

\src{snippet05}
左Kan扩张是一个余端：
$$\Lan_{K}D a = \int^i \cat{A}(K i, a)\times{}D i$$
因此它对应存在量词。形式化表示：

\begin{snip}{haskell}
Lan k d a = exists i. (k i -> a, d i)
\end{snip}
这可以通过GADT或全称量化数据构造器在Haskell中编码：

\src{snippet06}
对该数据结构的解释是：它包含一个函数，该函数接受某个未指定类型的\code{i}容器并生成\code{a}。同时它还有一个\code{i}容器。由于你不知道\code{i}的具体类型，唯一能做的操作就是取出\code{i}容器，通过自然变换将其重新打包为由函子\code{k}定义的容器，然后调用函数获取\code{a}。例如，若\d是树，\k是列表，你可以序列化该树，用生成的列表调用函数从而获得\code{a}。

左Kan扩张可用于计算函子的右伴随。我们知道乘积函子的右伴随是指数对象，让我们尝试用Kan扩张实现它：

\src{snippet07}
这确实与函数类型同构，见证如下函数对：

\src{snippet08}
注意，正如前面一般情况所述，我们执行了以下步骤：

\begin{enumerate}
  \tightlist
  \item
        取出\code{x}容器（此处只是平凡的恒等容器）和函数\code{f}。
  \item
        通过恒等函子与配对函子之间的自然变换重新打包容器。
  \item
        调用函数\code{f}。
\end{enumerate}

\section{自由函子}

Kan扩张的一个有趣应用是自由函子的构造。这是对以下实际问题的解答：
假设你有一个类型构造器——即对象之间的映射。能否基于该类型构造器定义一个函子？
换句话说，我们能否定义态射的映射，使得该类型构造器扩展为完整的自函子？

关键观察在于：类型构造器可以描述为以离散范畴为定义域的函子。离散范畴除了恒等态射外
没有其他态射。给定范畴 $\cat{C}$，我们总能通过丢弃所有非恒等态射构造出离散范畴 $\cat{|C|}$。
从 $\cat{|C|}$ 到 $\cat{C}$ 的函子 $F$ 就是一个简单的对象映射，这正是我们在Haskell中称为类型构造器的概念。
同时存在典范函子 $J$ 将 $\cat{|C|}$ 嵌入到 $\cat{C}$ 中：它在对象上（及恒等态射上）保持恒等。
若沿 $J$ 对 $F$ 作左Kan扩张，则得到从 $\cat{C}$ 到 $\cat{C}$ 的函子：
$$\Lan_{J}F a = \int^i \cat{C}(J i, a)\times{}F i$$
这被称为基于 $F$ 的自由函子。

在Haskell中可写作：

\src{snippet09}
事实上，对于任意类型构造器 \code{f}，\code{FreeF f} 都是函子：

\src{snippet10}
如你所见，自由函子通过记录函数及其参数来模拟函数提升。它通过记录函数组合来累积提升后的函数，
从而自动满足函子公理。该构造被用于论文
\urlref{http://okmij.org/ftp/Haskell/extensible/more.pdf}{Freer Monads,
  More Extensible Effects}。

同理，我们也可以使用右Kan扩张达到相同目的：

\src{snippet11}
容易验证这确实满足函子条件：

\src{snippet12}